\chapter{Supplementary Evidence and Reproducibility}

\section{Selected robustness results}
\label{app:robustness}

\begin{table}[H]
\centering
\caption{Robustness of the composite contemporaneous association.}
\label{tab:assocrobust}
\small
\begin{tabular}{lrrrr}
\toprule
Estimator & Coefficient & SE & 95\% CI & $p$ \\
\midrule
TWFE, country cluster & -0.0715 & 0.0353 & $[-0.1408,-0.0023]$ & 0.043 \\
TWFE, country + year cluster & -0.0715 & 0.0348 & $[-0.1398,-0.0033]$ & 0.040 \\
TWFE, Driscoll--Kraay & -0.0715 & 0.0152 & $[-0.1013,-0.0418]$ & $<0.001$ \\
First difference & -0.0243 & 0.0187 & $[-0.0610,0.0124]$ & 0.195 \\
\bottomrule
\end{tabular}
\end{table}

\begin{table}[H]
\centering
\caption{Primary temporal tests for composite corruption.}
\label{tab:directionappendix}
\small
\begin{tabularx}{\textwidth}{Yrrrr}
\toprule
Direction and test & Estimate / statistic & $p$ & Observations & Countries \\
\midrule
Populism $\rightarrow$ corruption, lead-lag L1 & $\theta_1=-0.056$ & 0.065 & 3,845 & 95 \\
Populism $\rightarrow$ corruption, lead-lag L2 lag 1 & $\theta_1=-0.048$ & 0.031 & 3,750 & 95 \\
Populism $\rightarrow$ corruption, Granger L2 joint & $\chi^2=0.753$ & 0.686 & 3,750 & 95 \\
Corruption $\rightarrow$ populism, Granger L2 joint & $\chi^2=2.848$ & 0.241 & 3,750 & 95 \\
Future populism placebo, joint & -- & 0.096 & 3,750 & 95 \\
Future corruption placebo, joint & -- & 0.841 & 3,750 & 95 \\
Populist-entry post-event joint & -- & 0.228 & 3,940 & 95 \\
Populist-entry pre-trend joint & -- & 0.983 & 3,940 & 95 \\
\bottomrule
\end{tabularx}
\end{table}

\section{Within-country information by corruption dimension}

\begin{table}[H]
\centering
\caption{Coverage, within variation, and persistence of corruption outcomes.}
\label{tab:dimensiondiagnostics}
\small
\begin{tabular}{lrrrrr}
\toprule
Outcome & $N$ & Missing & Within variance & Lag-1 corr. & Invariant countries \\
\midrule
Composite & 3,965 & 0 & 12.5\% & 0.993 & 2 \\
Executive & 3,965 & 0 & 17.2\% & 0.987 & 3 \\
Public sector & 3,965 & 0 & 16.0\% & 0.990 & 3 \\
Legislative & 3,770 & 195 & 12.4\% & 0.992 & 6 \\
Judicial & 3,965 & 0 & 11.1\% & 0.992 & 5 \\
\bottomrule
\end{tabular}
\end{table}

\section{Additional heterogeneity diagnostic}

\begin{figure}[H]
  \centering
  \includegraphics[width=0.78\textwidth]{democracy_lags.png}
  \caption{Populism-to-corruption lag estimates in lower- and higher-democracy country groups. The
  pattern suggests more negative lag-1 coefficients at lower democracy, but marginal effects are not
  statistically clear.}
  \label{fig:democracylags}
\end{figure}

\section{Reproducibility workflow}
\label{app:reproducibility}

The repository uses Python 3.12/3.13-compatible code managed by \codevar{uv}. Raw V-Dem and V-Party
files are imported to DuckDB, while the strict processed panel is committed so all reported analyses
can be reproduced without redistributing every raw file. From the repository root:

\begin{verbatim}
uv sync --extra dev --extra analysis
uv run python src/preprocessing.py
uv run pytest
uv run python scripts/generate_artifacts.py
uv run python scripts/generate_decomposition_artifacts.py
uv run python scripts/generate_report_assets.py
\end{verbatim}

The explanatory notebooks follow the analytical sequence:

\begin{enumerate}
  \item \codevar{01\_eda\_populism\_corruption.ipynb}: coverage, distributions, group contrasts,
  and country cases;
  \item \codevar{02\_preprocessing.ipynb}: party aggregation, coder filtering, forward-fill, join,
  and validation;
  \item \codevar{03\_association.ipynb}: between/within comparison, TWFE, interaction, machine
  learning, and robustness;
  \item \codevar{04\_causal\_direction.ipynb}: lead-lag, Granger-style tests, placebos, event studies,
  first differences, and electoral robustness;
  \item \codevar{05\_corruption\_decomposition.ipynb}: dimension-specific association, temporal
  direction, FDR correction, and measurement checks.
\end{enumerate}

The automated suite contains 33 passing tests. It checks preprocessing invariants, country-bounded
lags, duplicate keys, fixed sample construction, outcome orientation, synthetic signal recovery,
joint tests, multiplicity correction, and robustness coverage. All report figures are generated from
the committed panel or copied from reproducible analysis artifacts; no numerical result is transcribed
from an unversioned spreadsheet.

\section{Interpretation guide}

\begin{table}[H]
\centering
\caption{What each reported design can and cannot establish.}
\small
\begin{tabularx}{\textwidth}{>{\raggedright\arraybackslash}p{3.2cm}YY}
\toprule
Design & Supports & Does not establish \\
\midrule
Pooled / between OLS & Cross-country descriptive association & Within-country effect or direction \\
Two-way fixed effects & Same-year within-country co-movement after absorbed effects & Causal order or removal of time-varying confounding \\
Lead-lag FE & Temporal precedence in observed annual data & Incremental prediction beyond own persistence \\
Granger-style FE & Incremental predictive content of lagged proposed cause & Structural or intervention-based causality \\
Event study & Dynamic pattern around defined entries/shocks & Exogenous treatment assignment \\
First differences & Association among annual changes & Long-run level relationship \\
XGBoost + SHAP & Held-out-country prediction and fitted feature use & Statistical significance or causal importance \\
\bottomrule
\end{tabularx}
\end{table}
